% !TeX root = ../report_project_2stage_OTA.tex

\section{First-Order Design}

The OTA is first sized analytically to obtain a coherent initial operating
point before transistor-level optimization. These calculations are intended
as first-order design guidance rather than exact predictions of the final
BSIM-based circuit.

\subsection{Bias-Current Selection}

The diode-connected PMOS \texttt{PTAIL\_0} defines the reference bias
current

\[
I = 500~\mathrm{nA}.
\]

The first-stage PMOS current source is initially mirrored with a ratio of
one:

\[
I_{\mathrm{PTAIL\_1}} = I = 500~\mathrm{nA}.
\]

At zero differential input, the differential-pair current therefore splits
equally:

\[
I_{\mathrm{PDIFF\_P}}
=
I_{\mathrm{PDIFF\_N}}
=
\frac{I}{2}
=
250~\mathrm{nA}.
\]

The second-stage bias current is expressed as a multiple of the reference
current:

\[
I_{\mathrm{PTAIL\_2}} = N I.
\]

The value of $N$ is derived from the required second-stage transconductance
and stability constraints.


\subsection{Input-Stage Transconductance}

For the low-power input differential pair, an initial transconductance
efficiency of

\[
\frac{g_m}{I_D}
\approx
20~\mathrm{V}^{-1}
\]

is selected.

With

\[
I_D = 250~\mathrm{nA},
\]

the expected input-device transconductance is

\[
g_{m1}
=
\frac{g_m}{I_D}I_D
=
20 \times 250~\mathrm{nA}
=
5~\mu\mathrm{S}.
\]


\subsection{Miller Compensation Capacitor}

For a classical Miller-compensated two-stage OTA, the unity-gain bandwidth
is approximated by

\[
GBW
\approx
\frac{g_{m1}}
     {2\pi C_{\mathrm{MILLER}}}.
\]

Using the target

\[
GBW = 1~\mathrm{MHz},
\]

the required compensation capacitance is

\[
C_{\mathrm{MILLER}}
=
\frac{5~\mu\mathrm{S}}
     {2\pi \times 1~\mathrm{MHz}}
\approx
0.80~\mathrm{pF}.
\]

The initial compensation value is therefore

\[
\boxed{
C_{\mathrm{MILLER}}
\approx
0.8~\mathrm{pF}
}.
\]

A SKY130A MIM capacitor is used for the physical implementation of this
component.


\subsection{Second-Stage Transconductance}

The second-stage transconductance must place the non-dominant output pole
sufficiently above the unity-gain frequency.

Using the first-order approximation

\[
\omega_{p2}
\approx
\frac{g_{m2}}{C_L}
\]

and targeting

\[
\omega_{p2}
\gtrsim
2.2\,\omega_u
\]

for approximately $60^\circ$ phase margin gives

\[
g_{m2}
\gtrsim
2.2\,
g_{m1}
\frac{C_L}
     {C_{\mathrm{MILLER}}}.
\]

For the worst-case initial load

\[
C_L = 4~\mathrm{pF},
\]

the required second-stage transconductance is approximately

\[
g_{m2}
\gtrsim
2.2
\times
5~\mu\mathrm{S}
\times
\frac{4~\mathrm{pF}}
     {0.8~\mathrm{pF}}
=
55~\mu\mathrm{S}.
\]

The initial target is therefore

\[
\boxed{
g_{m2}
\approx
55\text{--}60~\mu\mathrm{S}
}.
\]


\subsection{Second-Stage Bias Ratio}

An initial second-stage current of

\[
4I = 2~\mu\mathrm{A}
\]

would require

\[
\frac{g_m}{I_D}
\approx
\frac{55~\mu\mathrm{S}}
     {2~\mu\mathrm{A}}
=
27.5~\mathrm{V}^{-1}.
\]

This operating point is unnecessarily aggressive for the initial design.

Increasing the mirror ratio to

\[
N = 6
\]

gives

\[
I_{\mathrm{PTAIL\_2}}
=
6I
=
3~\mu\mathrm{A}.
\]

The corresponding required transconductance efficiency becomes

\[
\frac{g_m}{I_D}
\approx
\frac{55~\mu\mathrm{S}}
     {3~\mu\mathrm{A}}
\approx
18.3~\mathrm{V}^{-1},
\]

which is a more suitable initial operating region.

The initial bias-current ratio is therefore selected as

\[
\boxed{
I_{\mathrm{PTAIL\_0}}
:
I_{\mathrm{PTAIL\_1}}
:
I_{\mathrm{PTAIL\_2}}
=
1:1:6
}.
\]


\subsection{First-Order Slew-Rate Estimate}

Although slew rate is not a primary design objective, it provides a useful
large-signal sanity check.

Assuming that the first-stage bias current limits the charging or
discharging of the Miller capacitor,

\[
SR
\approx
\frac{I_{\mathrm{PTAIL\_1}}}
     {C_{\mathrm{MILLER}}}.
\]

Therefore,

\[
SR
\approx
\frac{500~\mathrm{nA}}
     {0.8~\mathrm{pF}}
\approx
0.63~\mathrm{V/\mu s}.
\]

This value is considered sufficient as an initial estimate for the intended
low-power precision application.

\subsection{Initial Transistor Sizing}

The transistor dimensions are obtained from the SKY130A device
characterization performed in the preceding project.

For each device group, the drain current, target $g_m/I_D$ and channel
length are first selected from the circuit-level requirements. The
characterized current density $I_D/W$ is then used to determine the
required width.

\subsubsection{Input Differential Pair}

The PMOS differential pair carries

\[
I_D = 250~\mathrm{nA}
\]

per device and targets

\[
\frac{g_m}{I_D} \approx 20~\mathrm{V}^{-1}.
\]

A channel-length comparison showed that increasing the length from
$1~\mu\mathrm{m}$ to $2~\mu\mathrm{m}$ strongly increased device area
for comparatively limited circuit-level benefit.

The initial geometry is therefore selected as

\[
\boxed{
W_{\mathrm{PDIFF}} = 89~\mu\mathrm{m},
\qquad
L_{\mathrm{PDIFF}} = 1~\mu\mathrm{m}
}.
\]

At this operating point, the characterized intrinsic gain is approximately

\[
\frac{g_m}{g_{ds}} \approx 478.
\]


\subsubsection{First-Stage NMOS Mirror Load}

The two NMOS mirror devices also carry

\[
I_D = 250~\mathrm{nA}.
\]

A preliminary target of $g_m/I_D \approx 15~\mathrm{V}^{-1}$ resulted in
a calculated width below the minimum practical device width.

The width was therefore fixed to

\[
W = 0.42~\mu\mathrm{m},
\]

and the corresponding operating point was re-evaluated from the
characterized data.

For $L=1~\mu\mathrm{m}$, the resulting values are approximately

\[
\frac{g_m}{I_D} \approx 19.5~\mathrm{V}^{-1}
\]

and

\[
\frac{g_m}{g_{ds}} \approx 210.
\]

The initial geometry is therefore

\[
\boxed{
W_{\mathrm{NLOAD0,1}} = 0.42~\mu\mathrm{m},
\qquad
L_{\mathrm{NLOAD0,1}} = 1~\mu\mathrm{m}
}.
\]


\subsubsection{Second-Stage Gain Device}

The NMOS second-stage gain transistor carries

\[
I_D = 3~\mu\mathrm{A}
\]

and targets

\[
\frac{g_m}{I_D} \approx 18.3~\mathrm{V}^{-1}.
\]

The characterized data showed little intrinsic-gain improvement between
$L=1~\mu\mathrm{m}$ and $L=2~\mu\mathrm{m}$, while the required width
approximately doubled.

The selected initial geometry is therefore

\[
\boxed{
W_{\mathrm{NLOAD2}} = 3.58~\mu\mathrm{m},
\qquad
L_{\mathrm{NLOAD2}} = 1~\mu\mathrm{m}
}.
\]


\subsubsection{PMOS Bias Mirror}

The PMOS bias devices are operated around

\[
\frac{g_m}{I_D} \approx 15~\mathrm{V}^{-1}.
\]

For the reference current of $500~\mathrm{nA}$, the selected geometry is

\[
W_{\mathrm{PTAIL0}}
=
W_{\mathrm{PTAIL1}}
=
1.76~\mu\mathrm{m},
\qquad
L = 1~\mu\mathrm{m}.
\]

The second-stage current source uses the same operating point and a
$6:1$ mirror ratio:

\[
W_{\mathrm{PTAIL2}}
=
6 \times 1.76
=
10.56~\mu\mathrm{m}.
\]


\subsubsection{Initial Sizing Summary}

\begin{table}[ht]
    \centering
    \caption{Initial transistor sizing obtained from first-order design and
    SKY130A characterized device data.}
    \label{tab:initial_sizing}
    \begin{tabular}{lcccc}
        \hline
        Device & $I_D$ & $g_m/I_D$ & $W$ ($\mu$m) & $L$ ($\mu$m) \\
        \hline
        \texttt{PTAIL\_0} & 500 nA & 14.9 & 1.76 & 1.0 \\
        \texttt{PTAIL\_1} & 500 nA & 14.9 & 1.76 & 1.0 \\
        \texttt{PDIFF\_N/P} & 250 nA & 20.0 & 89.0 & 1.0 \\
        \texttt{NLOAD\_0/1} & 250 nA & 19.5 & 0.42 & 1.0 \\
        \texttt{PTAIL\_2} & 3 $\mu$A & 14.9 & 10.56 & 1.0 \\
        \texttt{NLOAD\_2} & 3 $\mu$A & 18.2 & 3.58 & 1.0 \\
        \hline
    \end{tabular}
\end{table}


\subsection{First-Order Gain Estimate}

Using the characterized output conductances at the selected operating
points, the first-stage gain is estimated as

\[
A_{v1}
\approx
g_{m1}
\left(
r_{o,\mathrm{PDIFF}}
\parallel
r_{o,\mathrm{NLOAD1}}
\right)
\approx
43.5~\mathrm{dB}.
\]

For the second stage,

\[
A_{v2}
\approx
g_{m2}
\left(
r_{o,\mathrm{NLOAD2}}
\parallel
r_{o,\mathrm{PTAIL2}}
\right)
\approx
43.7~\mathrm{dB}.
\]

The expected total low-frequency gain is therefore

\[
\boxed{
A_0
\approx
A_{v1}+A_{v2}
\approx
87~\mathrm{dB}
}.
\]

This provides approximately $7$ dB of nominal first-order margin relative
to the $80$ dB design target. The result will subsequently be verified
using transistor-level simulation.
